\documentclass[12pt]{article}

\include{../style}

\begin{document}

\title{Cosc 30: Discrete Mathematics}

\author{Prishita Dharampal}
\date{}
\maketitle
\vspace{-0.3in}


\begin{problem}{1}

    Consider the following table on the effect of a drug treatment on male and female patients.
    If we associate $C$ (connoting cause) with taking a certain drug and $E$ (connoting effect)
    with recovery, then the drug seems to be harmful to both males and females. A new patient
    comes in; do we use the drug or do we not, if our decision is purely based on whether the
    drug increases the chance of recovery among the whole sample population?

    (To think about later: the data seems to suggest when we further know the gender of the
    patient, we should not use the treatment. What gives?)

    \begin{center}
    \begin{tabular}{lcccc}
    \multicolumn{5}{c}{} \\
    \textbf{Males} & $E$ & $\neg E$ & Total & Recovery Rate \\
    \hline
    Drug ($C$) & 18 & 12 & 30 & 60\% \\
    No drug ($\neg C$) & 7 & 3 & 10 & 70\% \\
    \hline
    Total & 25 & 15 & 40 & \\
    \hline
    \end{tabular}

    \vspace{1em}

    \begin{tabular}{lcccc}
    \multicolumn{5}{c}{}\\
    \textbf{Females} & $E$ & $\neg E$ & Total & Recovery Rate \\
    \hline
    Drug ($C$) & 2 & 8 & 10 & 20\% \\
    No drug ($\neg C$) & 9 & 21 & 30 & 30\% \\
    \hline
    Total & 11 & 29 & 40 & \\
    \hline
    \end{tabular}
    \end{center}
    Figure 1. Recovery rates under treatment ($C$) and control ($\neg C$) for males and females.
\end{problem}

\begin{solution}
    \bbni

    Data: 
    \begin{itemize}
        \item With Drug: $18/30$ male and $2/10$ female patients recovered
        
        \[P(E|C) = \frac{18+2}{30+10} = 0.5\]
        
        \item Without Drug: $7/10$ male and $9/30$ female patients recovered
        
        \[P(E|\neg C) = \frac{7+9}{40} = 0.4\]
    \end{itemize}
    Based purely on the probability of recovery in the whole population sample, the patient should be given the drug. 

    However if we look at the gender samples separately, the recovery rates look like: 

    for men: 
    \[P(E|C) = \frac{18}{30} = 0.6 \qquad P(E|\neg C) = \frac{7}{10} = 0.7\]

    for women: 
    \[P(E|C) = \frac{2}{10} = 0.2 \qquad P(E|\neg C) = \frac{9}{30} = 0.3\]

    that is, patients have a higher chance of recovery if they don't take the drug. 
\end{solution}

\newpage 

\begin{problem}{2}

    There are two dice on the table. One of them is a fair dice, and the other is a loaded dice
    that always lands on six. You cannot distinguish the two dice from their appearances.

    Suppose you choose one of the dice at random, throw the dice $k$ times for some number
    $k$, and each time the dice lands on six. How large does $k$ have to be before you are
    99\% sure the dice you chose is the loaded dice?

    (To think about later: You must have thrown that many sixes in a row before. What gives?)
\end{problem}

\begin{solution}
    \bbni

    Let the event of choosing fair dice be $F$ and the loaded dice be $L$, then the probability of picking either dice is: 
    \[P( F) = P( L) = \frac{1}{2} = 0.5\]

    The probability $P(R_k)$ of rolling $k$ sixes in a row given a certain dice is, 
    \[P(R_k |L ) = 1, \qquad P(R_k|F) = \frac{1}{6^k}\]

    Then, to be $99\%$ certain that the dice we chose is the loaded dice after $k$ tries we calculate $P(L|R_k)$
    \begin{align*}
        P(L|R_k) &= \frac{P(R_k|L) \cdot P(L)}{P(R_k|L) \cdot P(L) + P(R_k|F) \cdot P(F)} \\ 
        &= \frac{1 \cdot \fr{1}{2}}{1 \cdot \fr{1}{2} + \fr{1}{6^k} \cdot \fr{1}{2}} \\ 
        &= \frac{1}{1 + \fr{1}{6^k}} \\ 
    \end{align*}
    and we want, 
    \begin{align*}
        P(L|R_k) = \frac{1}{1 + \fr{1}{6^k}} &\geq 0.99  \\ 
        \log(1/0.99 -1) &\geq k \log(1/6)\\ \\
    \intertext{Approximating $\log(1/0.99 -1)  \approx -2, \, \log(1/6) \approx -0.77$} \\ 
        -2/-0.77 =2.5 &\leq k  
    \end{align*}
    We get $k=3$. If we roll $3$ sixes in a row before we are $99\%$ sure that the chosen dice is loaded.
\end{solution}

\end{document}